\SourceLeaf{396}%
und
\[
\tfrac{c+d+e}{3}-\tfrac13\BareRoot{(cc+dd+ee+cd+ce+de)}
\]
sind.

Die Wurzelgrösse \(\BareRoot{(cc+dd+ee+cd+ce+de)}\) lässt sich auch in die Form setzen
\[
\BareRoot{[(d-c)^2+(d-c)(e-d)+(e-d)^2]}
\]
sie ist folglich die dritte Seite eines Dreiecks, in welchem zwei Seiten \(d-c\) und \(e-d\) sind, und der eingeschlossene Winkel \(=120^0\). Beschreibt man also über \(CD\) ein gleichseitiges Dreieck, dessen Spitze \(F\), so ist \(EF\) jener Wurzelgrösse gleich, wonach sich also die beiden Werthe von \(r\) leicht construiren lassen. Man kann leicht zeigen, dass der eine dieser Werthe zwischen \(c\) und \(d\), der andere zwischen \(d\) und \(e\) fallen muss, und dass nur dem erstern der Mittelpunkt der grössten Ellipse wirklich entspricht; für den andern wird nemlich
\[
\gamma+\delta r+\varepsilon rr+\zeta r^3
\]
nicht ein grösstes, sondern ein kleinstes werden, oder vielmehr den grössten \textit{negativen} Werth erhalten, dem also nur ein imaginärer Werth von \(\alpha\beta\) entsprechen kann. Man sieht leicht, dass dieser sich auf eine Hyperbel beziehen muss.

Sobald übrigens der Mittelpunkt der verlangten Ellipse gefunden ist, hat die Bestimmung der übrigen unbekannten Grössen keine Schwierigkeit. Aus \(\theta\) und \(r\) findet man \(t\); aus \(t\) und \(u\) dann ferner \(\alpha\) und \(\beta\), und dann aus einer oder einigen der obigen Gleichungen \(\psi\). Dadurch sind also sowohl die Dimensionen der Ellipse, als ihre Lage vollkommen bestimmt.

Ich muss übrigens noch bemerken, dass das hier aufgelöste Problem mit dem neulich in der \textit{Monatl. Corresp.} aufgegebenen nicht ganz einerlei ist. Es gibt nemlich Fälle, wo die grösste \textit{innerhalb} eines Vierecks zu beschreibende Ellipse eine der vier Seiten des Vierecks nicht berührt. Die nähere Betrachtung dieser Fälle gehört aber hier nicht zu meiner Absicht.

\par\medskip\Needspace{4\baselineskip}
\pdfbookmark[2]{Directer Beweis des obigen Theorems die Vierecke betreffend.}{ellipse-directer-beweis}
\begin{center}\textit{Directer Beweis des obigen Theorems die Vierecke betreffend.}\end{center}
\hspace*{\parindent}Es seien \(A,B,C,D\) die vier Winkelpunkte des Vierecks; \(E\) der Durchschnitt von \(AB\) und \(DC\); \(F\) der Durchschnitt von \(BC\) und \(AD\); \(G,H\) und \(I\)
